\documentclass[11pt,a4paper]{article}
\usepackage[T1]{fontenc}
\usepackage{lmodern,array,amsmath,amssymb,booktabs,enumitem,fancyhdr,lastpage,needspace,microtype}
\usepackage[margin=19mm,headheight=14pt,footskip=11mm]{geometry}
\usepackage[hidelinks]{hyperref}
\setlength{\parindent}{0pt}
\setlength{\parskip}{5pt}
\setlength{\emergencystretch}{2em}
\pagestyle{fancy}\fancyhf{}
\fancyhead[L]{\small CS182 --- Lecture 05}
\fancyhead[R]{\small Solutions}
\fancyfoot[L]{\small Independent practice quiz}
\fancyfoot[R]{\small Page \thepage\ of \pageref{LastPage}}
\newcommand{\SetAnswer}[2]{\expandafter\gdef\csname QuizAnswer#1\endcsname{#2}}
\newcommand{\GetAnswer}[1]{\csname QuizAnswer#1\endcsname}
\newlength{\BlankAnswerHeight}\setlength{\BlankAnswerHeight}{16mm}
\newcommand{\AnswerBlank}{\leavevmode\par\vspace{\BlankAnswerHeight}}

\begin{document}
\begin{center}{\Large\bfseries CS182: Introduction to Machine Learning}\par{\Large\bfseries Lecture 05: Bayesian Decision Theory}\par{\large Worked Solutions and Marking Guide}\end{center}
\textbf{Scope:} Bayesian\_Decision\_Theory.pdf. Independent practice material based on the official 2026 Fall lecture.\par\textbf{Total:} 100 marks. \quad \textbf{Suggested time:} 120 minutes.
\begin{center}\begin{tabular}{lrrr}\toprule
Section & Questions & Each & Marks\\\midrule
A. Misconceptions: True/False & 12 & 1 & 12\\
B. Multiple response & 12 & 2 & 24\\
C. Proofs and formula derivations & 4 & 10 & 40\\
D. Applications & 2 & 12 & 24\\\midrule
Total & 30 & & 100\\\bottomrule\end{tabular}\end{center}
\textbf{Scoring.} In Section B each item has exactly four options, of which one,
two, or three are correct. The complete correct set earns 2 marks; a nonempty
proper subset earns 1 mark; any incorrect option or no selection earns 0.
T/F answers earn 1 mark only when correct. Show reasoning in Sections C and D;
part marks are available for valid intermediate work. Equivalent correct methods receive credit.
\par\medskip\textbf{Instructions.} Answer in English. Use natural logarithms.
Exact values and justified approximations are accepted. All vectors are columns
unless stated otherwise. Use local assumptions, label conventions, and tie rules.
\par\medskip\textbf{Quick key --- Section A}
\begin{center}\renewcommand{\arraystretch}{1.15}\begin{tabular}{|c|c|c|c|c|c|}
\hline
1 & 2 & 3 & 4 & 5 & 6\\\hline
\makebox[20mm]{\rule{0pt}{5mm}F} & \makebox[20mm]{\rule{0pt}{5mm}T} & \makebox[20mm]{\rule{0pt}{5mm}F} & \makebox[20mm]{\rule{0pt}{5mm}T} & \makebox[20mm]{\rule{0pt}{5mm}T} & \makebox[20mm]{\rule{0pt}{5mm}F}\\\hline
\end{tabular}\end{center}
\begin{center}\renewcommand{\arraystretch}{1.15}\begin{tabular}{|c|c|c|c|c|c|}
\hline
7 & 8 & 9 & 10 & 11 & 12\\\hline
\makebox[20mm]{\rule{0pt}{5mm}F} & \makebox[20mm]{\rule{0pt}{5mm}F} & \makebox[20mm]{\rule{0pt}{5mm}F} & \makebox[20mm]{\rule{0pt}{5mm}T} & \makebox[20mm]{\rule{0pt}{5mm}T} & \makebox[20mm]{\rule{0pt}{5mm}F}\\\hline
\end{tabular}\end{center}
\par\medskip\textbf{Quick key --- Section B}
\begin{center}\renewcommand{\arraystretch}{1.15}\begin{tabular}{|c|c|c|c|c|c|}
\hline
13 & 14 & 15 & 16 & 17 & 18\\\hline
\makebox[20mm]{\rule{0pt}{5mm}BCD} & \makebox[20mm]{\rule{0pt}{5mm}ACD} & \makebox[20mm]{\rule{0pt}{5mm}BCD} & \makebox[20mm]{\rule{0pt}{5mm}CD} & \makebox[20mm]{\rule{0pt}{5mm}ABC} & \makebox[20mm]{\rule{0pt}{5mm}ABD}\\\hline
\end{tabular}\end{center}
\begin{center}\renewcommand{\arraystretch}{1.15}\begin{tabular}{|c|c|c|c|c|c|}
\hline
19 & 20 & 21 & 22 & 23 & 24\\\hline
\makebox[20mm]{\rule{0pt}{5mm}A} & \makebox[20mm]{\rule{0pt}{5mm}AD} & \makebox[20mm]{\rule{0pt}{5mm}AC} & \makebox[20mm]{\rule{0pt}{5mm}ABC} & \makebox[20mm]{\rule{0pt}{5mm}BC} & \makebox[20mm]{\rule{0pt}{5mm}ABC}\\\hline
\end{tabular}\end{center}
\newpage\section*{Section A. True/False --- Explanations}
\par\noindent\begin{minipage}{\linewidth}
\textbf{1.} \textbf{False.} A continuous density can exceed 1; its integral, rather than its height, is normalized. Class posterior probabilities \(P(C_i\mid x)\) do lie in \([0,1]\).

{\footnotesize\textit{Source: PDF pp.11-12.}}
\end{minipage}\par\vspace{3mm}
\par\noindent\begin{minipage}{\linewidth}
\textbf{2.} \textbf{True.} For example, if \(P(C_1\mid x)=0.3\), the loss of predicting \(C_1\) when \(C_2\) is true is 1, and the opposite error costs 4, then \(R(\alpha_1\mid x)=0.7<1.2=R(\alpha_2\mid x)\).

{\footnotesize\textit{Source: PDF p.15; supplementary numerical illustration.}}
\end{minipage}\par\vspace{3mm}
\par\noindent\begin{minipage}{\linewidth}
\textbf{3.} \textbf{False.} The tied scores must also be largest. If another class has a strictly larger score, the equality is hidden inside that class's decision region.

{\footnotesize\textit{Source: PDF p.23.}}
\end{minipage}\par\vspace{3mm}
\par\noindent\begin{minipage}{\linewidth}
\textbf{4.} \textbf{True.} Every pairwise difference is multiplied by the same positive number \(a(x)\). Its sign and any zero difference are preserved pointwise.

{\footnotesize\textit{Source: PDF p.21; supplementary pointwise extension.}}
\end{minipage}\par\vspace{3mm}
\par\noindent\begin{minipage}{\linewidth}
\textbf{5.} \textbf{True.} The error of choosing \(C_i\) is \(1-P(C_i)\); minimizing it is equivalent to maximizing the prior. Ties allow several optimal constant decisions.

{\footnotesize\textit{Source: PDF pp.6-7,16.}}
\end{minipage}\par\vspace{3mm}
\par\noindent\begin{minipage}{\linewidth}
\textbf{6.} \textbf{False.} Within-class monotonicity does not preserve comparisons across classes. Scores \(g_1=2,g_2=1\) choose \(C_1\), whereas \(f_1(t)=t\) and \(f_2(t)=t+10\) choose \(C_2\). The same increasing function must be used for all classes.

{\footnotesize\textit{Source: PDF p.21.}}
\end{minipage}\par\vspace{3mm}
\par\noindent\begin{minipage}{\linewidth}
\textbf{7.} \textbf{False.} For a fixed \(x\), the same positive denominator does not change the maximizing class. It is still required to make \(\sum_iP(C_i\mid x)=1\).

{\footnotesize\textit{Source: PDF pp.11-12,21.}}
\end{minipage}\par\vspace{3mm}
\par\noindent\begin{minipage}{\linewidth}
\textbf{8.} \textbf{False.} The relevant scores are \(p(x\mid C_i)P(C_i)\). Maximum likelihood classification agrees automatically with MAP when the priors are equal; unequal priors can reverse the ordering.

{\footnotesize\textit{Source: PDF pp.11-12,16.}}
\end{minipage}\par\vspace{3mm}
\par\noindent\begin{minipage}{\linewidth}
\textbf{9.} \textbf{False.} The mixed risk is a weighted average of the action risks and cannot be below their minimum. Randomization among minimizers can tie that minimum.

{\footnotesize\textit{Source: PDF p.15; supplementary randomization argument.}}
\end{minipage}\par\vspace{3mm}
\par\noindent\begin{minipage}{\linewidth}
\textbf{10.} \textbf{True.} Their risks are both \(1-\max_kP(C_k\mid x)\). A tie-breaking convention determines the output label without changing this minimum risk.

{\footnotesize\textit{Source: PDF pp.12,16.}}
\end{minipage}\par\vspace{3mm}
\par\noindent\begin{minipage}{\linewidth}
\textbf{11.} \textbf{True.} The best classification risk is \(1-\max_iP(C_i\mid x)\), while rejection has risk \(\lambda\). The strict comparison gives the stated threshold; equality is a risk tie.

{\footnotesize\textit{Source: PDF pp.17-18.}}
\end{minipage}\par\vspace{3mm}
\par\noindent\begin{minipage}{\linewidth}
\textbf{12.} \textbf{False.} The required factorization is \(p(x\mid C_i)=\prod_jp(x_j\mid C_i)\), separately within each class. Mixing classes can produce dependent features even if this conditional factorization holds.

{\footnotesize\textit{Source: PDF p.13.}}
\end{minipage}\par\vspace{3mm}
\newpage\section*{Section B. Multiple Response --- Explanations}
\par\noindent\begin{minipage}{\linewidth}
% BEGIN QUESTION Q13 / 题目与选项区
\textbf{13.} \textbf{Correct set: B, C, D.}
\begin{enumerate}[label=(\Alph*),leftmargin=8mm,itemsep=2pt,topsep=2pt]
\item \textbf{Incorrect.} Its score is always strictly smaller than \(g_3\).
\item \textbf{Correct.} It is \(x<-1/\sqrt2\) or \(x>1/\sqrt2\).
\item \textbf{Correct.} These solve \(g_1=g_3=1/2\), and those tied scores exceed \(g_2=0\).
\item \textbf{Correct.} At both points \(g_3=1/2\) exceeds the tied value 0.
\end{enumerate}
{\footnotesize\textit{Source: PDF p.23; supplementary algebraic example.}}
% END QUESTION Q13

\end{minipage}\par\vspace{3mm}
\par\noindent\begin{minipage}{\linewidth}
% BEGIN QUESTION Q14 / 题目与选项区
\textbf{14.} \textbf{Correct set: A, C, D.}
\begin{enumerate}[label=(\Alph*),leftmargin=8mm,itemsep=2pt,topsep=2pt]
\item \textbf{Correct.} The constant rule is available among all rules using \(X\), so optimization cannot have a larger risk.
\item \textbf{Incorrect.} A feature-dependent rule may ignore useful posterior information or deliberately choose a low-posterior class.
\item \textbf{Correct.} The pointwise minimum error \(1-\max_iP(C_i\mid X)\) is zero almost surely.
\item \textbf{Correct.} Then \(E[1-\max_iP(C_i\mid X)]=1-\max_iP(C_i)\).
\end{enumerate}
{\footnotesize\textit{Source: PDF pp.7,9-12,15-16; supplementary comparison of rule classes.}}
% END QUESTION Q14

\end{minipage}\par\vspace{3mm}
\par\noindent\begin{minipage}{\linewidth}
% BEGIN QUESTION Q15 / 题目与选项区
\textbf{15.} \textbf{Correct set: B, C, D.}
\begin{enumerate}[label=(\Alph*),leftmargin=8mm,itemsep=2pt,topsep=2pt]
\item \textbf{Incorrect.} The acceptance threshold is \(1-0.55=0.45\), not \(0.55\).
\item \textbf{Correct.} The best class risk \(0.54\) already exceeds \(0.50\).
\item \textbf{Correct.} The best classification risk is \(1-0.46=0.54<0.55\).
\item \textbf{Correct.} Rejection is strictly optimal when \(1-\max_iP(C_i\mid x)>\lambda\); increasing \(\lambda\) makes this condition harder to satisfy.
\end{enumerate}
{\footnotesize\textit{Source: PDF pp.17-18.}}
% END QUESTION Q15

\end{minipage}\par\vspace{3mm}
\par\noindent\begin{minipage}{\linewidth}
% BEGIN QUESTION Q16 / 题目与选项区
\textbf{16.} \textbf{Correct set: C, D.}
\begin{enumerate}[label=(\Alph*),leftmargin=8mm,itemsep=2pt,topsep=2pt]
\item \textbf{Incorrect.} A negative multiplier reverses strict order, making maximizers into minimizers.
\item \textbf{Incorrect.} Class-dependent offsets can reverse cross-class comparisons.
\item \textbf{Correct.} Strict monotonicity preserves both strict order and ties across all classes.
\item \textbf{Correct.} The pairwise differences become \(a(x)[g_i(x)-g_j(x)]\), preserving their signs.
\end{enumerate}
{\footnotesize\textit{Source: PDF p.21; supplementary pointwise extension.}}
% END QUESTION Q16

\end{minipage}\par\vspace{3mm}
\newpage\section*{Section B. Multiple Response --- Explanations}
\par\noindent\begin{minipage}{\linewidth}
% BEGIN QUESTION Q17 / 题目与选项区
\textbf{17.} \textbf{Correct set: A, B, C.}
\begin{enumerate}[label=(\Alph*),leftmargin=8mm,itemsep=2pt,topsep=2pt]
\item \textbf{Correct.} Every action risk gains the same \(\sum_kb_kP(C_k\mid x)\).
\item \textbf{Correct.} Maximizing \(-R(\alpha_i\mid x)\) is equivalent to minimizing \(R(\alpha_i\mid x)\).
\item \textbf{Correct.} All action risks are multiplied by that constant, preserving order and ties.
\item \textbf{Incorrect.} Each action then gains a different offset, which can change the minimizing action.
\end{enumerate}
{\footnotesize\textit{Source: PDF pp.15,20; supplementary loss transformations.}}
% END QUESTION Q17

\end{minipage}\par\vspace{3mm}
\par\noindent\begin{minipage}{\linewidth}
% BEGIN QUESTION Q18 / 题目与选项区
\textbf{18.} \textbf{Correct set: A, B, D.}
\begin{enumerate}[label=(\Alph*),leftmargin=8mm,itemsep=2pt,topsep=2pt]
\item \textbf{Correct.} The log posterior odds equal \(\log L(x)+\log(1/3)\).
\item \textbf{Correct.} The joint-score comparison is \(L(x)>P(C_2)/P(C_1)=3\).
\item \textbf{Incorrect.} It is the same positive denominator for both posteriors and cancels.
\item \textbf{Correct.} The threshold is the prior odds \(0.9/0.1\).
\end{enumerate}
{\footnotesize\textit{Source: PDF pp.11-12,21-22.}}
% END QUESTION Q18

\end{minipage}\par\vspace{3mm}
\par\noindent\begin{minipage}{\linewidth}
% BEGIN QUESTION Q19 / 题目与选项区
\textbf{19.} \textbf{Correct set: A.}
\begin{enumerate}[label=(\Alph*),leftmargin=8mm,itemsep=2pt,topsep=2pt]
\item \textbf{Correct.} This is exactly the factorization \(p(x\mid C_i)=\prod_jp(x_j\mid C_i)\).
\item \textbf{Incorrect.} Conditioning can introduce dependence; independence in the mixture does not establish independence within each class.
\item \textbf{Incorrect.} Bayes scores remain \(P(C_i)\prod_jp(x_j\mid C_i)\), so unequal priors still matter.
\item \textbf{Incorrect.} Naive Bayes factors feature likelihoods, not posteriors of the class. Repeated posterior factors also repeatedly include the prior.
\end{enumerate}
{\footnotesize\textit{Source: PDF pp.11-13.}}
% END QUESTION Q19

\end{minipage}\par\vspace{3mm}
\par\noindent\begin{minipage}{\linewidth}
% BEGIN QUESTION Q20 / 题目与选项区
\textbf{20.} \textbf{Correct set: A, D.}
\begin{enumerate}[label=(\Alph*),leftmargin=8mm,itemsep=2pt,topsep=2pt]
\item \textbf{Correct.} The posterior vector is \((0.45,0.30,0.25)\), so \(C_1\) has the largest posterior.
\item \textbf{Incorrect.} The maximum posterior is \(0.45\); the error is its complement, \(0.55\).
\item \textbf{Incorrect.} It is \(0.18/0.40=0.45\); the likelihood alone is not the posterior.
\item \textbf{Correct.} Summing the joint scores gives \(0.18+0.12+0.10=0.40\).
\end{enumerate}
{\footnotesize\textit{Source: PDF pp.11-12,16.}}
% END QUESTION Q20

\end{minipage}\par\vspace{3mm}
\newpage\section*{Section B. Multiple Response --- Explanations}
\par\noindent\begin{minipage}{\linewidth}
% BEGIN QUESTION Q21 / 题目与选项区
\textbf{21.} \textbf{Correct set: A, C.}
\begin{enumerate}[label=(\Alph*),leftmargin=8mm,itemsep=2pt,topsep=2pt]
\item \textbf{Correct.} A convex combination of their equal risks remains \(0.6\).
\item \textbf{Incorrect.} Its risk is \(1-0.2=0.8\). The rejection threshold is \(1-0.6=0.4\), not \(0.6\).
\item \textbf{Correct.} Predicting either of the first two classes costs \(0.6\), as does rejection.
\item \textbf{Incorrect.} The two largest-posterior class actions tie with rejection. A reject-on-equality convention selects one of several minimizers.
\end{enumerate}
{\footnotesize\textit{Source: PDF pp.15-18; supplementary tie/randomization interpretation.}}
% END QUESTION Q21

\end{minipage}\par\vspace{3mm}
\par\noindent\begin{minipage}{\linewidth}
% BEGIN QUESTION Q22 / 题目与选项区
\textbf{22.} \textbf{Correct set: A, B, C.}
\begin{enumerate}[label=(\Alph*),leftmargin=8mm,itemsep=2pt,topsep=2pt]
\item \textbf{Correct.} This error occurs when \(C_2\) is true, giving \(8(1-0.55)=3.6\).
\item \textbf{Correct.} Its risk is \(2(0.55)=1.1<3.6\).
\item \textbf{Correct.} The inequality \(8(1-q)<2q\) is equivalent to \(q>0.8\).
\item \textbf{Incorrect.} Posterior maximization minimizes unit error loss; these asymmetric costs produce a different threshold.
\end{enumerate}
{\footnotesize\textit{Source: PDF pp.15-16; supplementary asymmetric-loss specialization.}}
% END QUESTION Q22

\end{minipage}\par\vspace{3mm}
\par\noindent\begin{minipage}{\linewidth}
% BEGIN QUESTION Q23 / 题目与选项区
\textbf{23.} \textbf{Correct set: B, C.}
\begin{enumerate}[label=(\Alph*),leftmargin=8mm,itemsep=2pt,topsep=2pt]
\item \textbf{Incorrect.} The mixed risk is \(0.2+0.3w_2\), which is at least \(0.2\).
\item \textbf{Correct.} Equality at the minimum requires \(w_2=0\), because its risk is strictly larger.
\item \textbf{Correct.} Its expected risk is \((0.2+0.2)/2=0.2\).
\item \textbf{Incorrect.} Any mixture supported on the first and third actions is optimal.
\end{enumerate}
{\footnotesize\textit{Source: PDF p.15; supplementary randomization derivation.}}
% END QUESTION Q23

\end{minipage}\par\vspace{3mm}
\par\noindent\begin{minipage}{\linewidth}
% BEGIN QUESTION Q24 / 题目与选项区
\textbf{24.} \textbf{Correct set: A, B, C.}
\begin{enumerate}[label=(\Alph*),leftmargin=8mm,itemsep=2pt,topsep=2pt]
\item \textbf{Correct.} It is \(4(0.1)+1(0.9)=1.3\); a density value may exceed 1.
\item \textbf{Correct.} It is \(0.4/1.3=4/13\).
\item \textbf{Correct.} Its joint score \(0.9\) exceeds \(0.4\), so its posterior is \(9/13\).
\item \textbf{Incorrect.} The prior-weighted likelihoods determine the posterior ordering.
\end{enumerate}
{\footnotesize\textit{Source: PDF pp.11-12,16.}}
% END QUESTION Q24

\end{minipage}\par\vspace{3mm}
\newpage\section*{Section C. Proofs and Formula Derivations}
\textbf{25. Why pointwise decisions minimize total risk}\hfill\textbf{[10 marks]}\par\medskip
\textbf{(a) [4 marks]} Conditioning first on \(X\) gives
\[
\mathcal R(\delta)=E_X\!\left[\sum_k\lambda_{\delta(X),k}P(C_k\mid X)\right]
=\int R(\delta(x)\mid x)p(x)\,dx.
\]
Use a sum instead of an integral for discrete \(X\). For a measurable minimizing selector \(\delta^*\), \(R(\delta^*(x)\mid x)\le R(\delta(x)\mid x)\) wherever posteriors are defined. Multiplication by \(p(x)\ge0\) and integration prove \(\mathcal R(\delta^*)\le\mathcal R(\delta)\). Choices on sets of \(X\)-probability zero do not affect total risk; ties can be resolved arbitrarily. Marking: conditional expectation (2), pointwise inequality and integration (1), ties/null sets (1).

\textbf{(b) [3 marks]} For action \(\alpha_i\),
\[
R(\alpha_i\mid x)=\sum_{k\ne i}P(C_k\mid x)=1-P(C_i\mid x).
\]
Thus any largest-posterior class is a conditional-risk minimizer. Substitution into the total-risk formula yields
\[
\mathcal R^*=E_X[1-\max_iP(C_i\mid X)].
\]
Marking: risk simplification (1), MAP minimizer (1), expected minimum error (1).

\textbf{(c) [3 marks]} Choose \(j\in\arg\max_iP(C_i)\). Pointwise \(\max_iP(C_i\mid X)\ge P(C_j\mid X)\). Taking expectations and using the law of total probability gives \(E[\max_iP(C_i\mid X)]\ge E[P(C_j\mid X)]=P(C_j)\). Subtracting from 1 proves the bound. The result concerns the true Bayes rule; an estimated or poorly chosen classifier need not minimize true risk. Marking: pointwise bound (1), expectation and final inequality (1), scope of guarantee (1).

\textit{Marking: use the displayed subpart allocations; award credit for the logical or numerical intermediate steps.}\par
{\footnotesize\textit{Source: PDF pp.15-16; supplementary integration and information comparison.}}
\newpage\section*{Section C. Proofs and Formula Derivations}
\textbf{26. Asymmetric errors and rejection thresholds}\hfill\textbf{[10 marks]}\par\medskip
\textbf{(a) [4 marks]} The risks are \(R_1=a(1-q)\) and \(R_2=bq\). Therefore \(R_1<R_2\) exactly when \(q>a/(a+b)\); reverse the inequality for \(C_2\), with a tie at equality. Equivalently \(q/(1-q)>a/b\). Bayes' rule gives
\[
\frac{q}{1-q}=\frac{p(x\mid C_1)P(C_1)}{p(x\mid C_2)P(C_2)},
\]
so predict \(C_1\) when
\[
\frac{p(x\mid C_1)}{p(x\mid C_2)}>
\frac{aP(C_2)}{bP(C_1)}.
\]
Marking: risks (1), posterior threshold and tie (1), posterior odds (1), likelihood-ratio threshold (1).

\textbf{(b) [3 marks]} Rejection is strictly optimal precisely when both \(r<a(1-q)\) and \(r<bq\), hence
\[
\frac rb<q<1-\frac ra.
\]
This interval has positive width exactly when \(r/b<1-r/a\), or \(r<ab/(a+b)\). At either endpoint rejection ties a class action; at \(r=ab/(a+b)\) it only ties at the original class threshold. For larger \(r\), rejection is never strictly optimal. Marking: two risk comparisons (1), interval (1), width condition and ties (1).

\textbf{(c) [3 marks]} The best class risk is \(1-M(x)\), where \(M(x)=\max_iP(C_i\mid x)\), and reject risk is \(\lambda\). Classification is strictly preferable exactly when \(M(x)>1-\lambda\); rejection is strictly preferable when \(M(x)<1-\lambda\). Since the \(K\) posteriors sum to 1, \(M(x)\ge1/K\). If \(\lambda>1-1/K\), then \(1-\lambda<1/K\le M(x)\), so a class action strictly beats rejection everywhere. At equality of the costs, the uniform posterior can produce a tie. Marking: risks (1), threshold/tie distinction (1), \(1/K\) bound and consequence (1).

\textit{Marking: use the displayed subpart allocations; award credit for the logical or numerical intermediate steps.}\par
{\footnotesize\textit{Source: PDF pp.15-18; supplementary asymmetric-cost extension.}}
\newpage\section*{Section C. Proofs and Formula Derivations}
\textbf{27. Equivalent discriminants and the role of strict monotonicity}\hfill\textbf{[10 marks]}\par\medskip
\textbf{(a) [4 marks]} For any two scores \(u,v\), strict monotonicity gives \(u>v\iff f(u)>f(v)\), and injectivity gives \(u=v\iff f(u)=f(v)\). Thus a class is at least as large as every competitor before transformation exactly when it is afterward. The full argmax set is unchanged. If \(f\) is merely nondecreasing, it can merge different scores: take \(g_1=2,g_2=1\) and \(f(t)=0\) for all \(t\). The original argmax is only \(C_1\), whereas both classes maximize transformed scores. Marking: order/tie argument (2), argmax conclusion (1), counterexample (1).

\textbf{(b) [3 marks]} Bayes' rule gives \(P(C_i\mid x)=p(x\mid C_i)P(C_i)/p(x)\). Multiplying all scores by common \(p(x)>0\) preserves their ordering and ties, so the posterior and joint scores have identical argmax sets. Applying the common strictly increasing logarithm yields \(\log p(x\mid C_i)+\log P(C_i)\), with the same maximizers. If a joint score is zero, the finite logarithm is undefined; one may compare the original scores or use the extended value \(\log0=-\infty\). At least one score is positive when \(p(x)>0\). Marking: common denominator (1), log decomposition (1), zero-score qualification (1).

\textbf{(c) [3 marks]} Subtract the two log joint scores:
\[
g(x)=\log\frac{p(x\mid C_1)}{p(x\mid C_2)}
+\log\frac{P(C_1)}{P(C_2)}.
\]
If \(g>0\), the posterior of \(C_1\) is larger; if \(g<0\), that of \(C_2\) is larger; \(g=0\) is a tie. For arbitrary losses, the optimal action minimizes \(\sum_k\lambda_{ik}P(C_k\mid x)\), which need not be the maximum-posterior class. Scores \(h_i(x)=-R(\alpha_i\mid x)\) give the correct maximizing action. Marking: log-ratio derivation (1), sign/tie rule (1), general-loss argument (1).

\textit{Marking: use the displayed subpart allocations; award credit for the logical or numerical intermediate steps.}\par
{\footnotesize\textit{Source: PDF pp.20-22; supplementary tie-preserving conditions.}}
\newpage\section*{Section C. Proofs and Formula Derivations}
\textbf{28. What a decision boundary actually represents}\hfill\textbf{[10 marks]}\par\medskip
\textbf{(a) [4 marks]} Each comparison is
\[
(w_i-w_j)^\mathsf Tx+(b_i-b_j)\ge0,
\]
a closed half-space (possibly the whole space or empty). Thus \(\mathcal R_i\) is an intersection of half-spaces. If \(u,v\) satisfy every inequality, then for \(0\le t\le1\) each left side at \(tu+(1-t)v\) is the corresponding convex combination of nonnegative values, so the segment also lies in \(\mathcal R_i\). Equality \(g_i=g_j\) is a hyperplane when \(w_i\ne w_j\), but only its portions where these scores are maximal can be decision boundaries. Marking: inequalities (1), intersection/segment proof (2), equality and maximality condition (1).

\textbf{(b) [3 marks]} Class \(C_1\) wins strictly when \(x^2>4\), that is, \(x<-2\) or \(x>2\). These are two disconnected components. Class \(C_2\) wins strictly when \(-2<x<2\). At \(x=\pm2\), the two scores tie for the maximum and form the boundaries. The weak regions include the endpoints, and any tie convention assigns them without affecting the disconnected strict regions. Marking: exterior region (1), interior and boundaries (1), disconnectedness/ties (1).

\textbf{(c) [3 marks]} The equality \(g_1=g_2\) gives \(x=\pm\sqrt2\), while \(g_1=g_3\) gives \(x=\pm1\); \(g_2=g_3\) has no solution. At \(\pm\sqrt2\), the tied first two scores are 0 but \(g_3=1\), so these points lie strictly in \(C_3\)'s region. At \(\pm1\), \(g_1=g_3=1>0=g_2\), and the winner changes between \(C_1\) and \(C_3\). The only boundaries are \(\pm1\). Marking: pairwise equalities (1), largest-score check (1), actual boundaries and counterexample conclusion (1).

\textit{Marking: use the displayed subpart allocations; award credit for the logical or numerical intermediate steps.}\par
{\footnotesize\textit{Source: PDF p.23; supplementary affine and quadratic examples.}}
\newpage\section*{Section D. Applications}
\textbf{29. A sensor with asymmetric mistakes and a review action}\hfill\textbf{[12 marks]}\par\medskip
\textbf{(a) [4 marks]} The two joint probabilities at \(A\) are \(0.175,0.065\); at \(B\), \(0.105,0.260\); at \(C\), \(0.070,0.325\). Therefore
\[
P(A)=0.240,\quad P(B)=0.365,\quad P(C)=0.395,
\]
which sum to 1. The posteriors are
\[
q_A=\frac{35}{48}\approx0.729167,\quad
q_B=\frac{21}{73}\approx0.287671,\quad
q_C=\frac{14}{79}\approx0.177215.
\]
Marking: joint scores (1), report probabilities (1), posterior values (2).

\textbf{(b) [4 marks]} For posterior \(q\), the risks are \(3(1-q)\), \(9q\), and \(1.7\). At \(A\), these are \(13/16=0.8125\), \(105/16=6.5625\), and \(1.7\), so predict \(C_1\). At \(B\), they are \(156/73\approx2.136986\), \(189/73\approx2.589041\), and \(1.7\), so review. At \(C\), they are \(195/79\approx2.468354\), \(126/79\approx1.594937\), and \(1.7\), so predict \(C_2\). Equivalently, review is preferred for \(17/90<q<13/30\). Marking: risk formulas (1), three risk triples with correct actions (3).

\textbf{(c) [4 marks]} The review-enabled expected loss is
\[
0.240(0.8125)+0.365(1.7)+0.395\left(\frac{126}{79}\right)
=0.195+0.6205+0.630=1.4455.
\]
Without review, the posterior threshold is \(3/(3+9)=0.25\), giving \(C_1\) at \(A,B\) and \(C_2\) at \(C\). Its loss is \(3(0.065)+3(0.260)+9(0.070)=1.605\). MAP uses threshold \(0.5\), giving \(C_1\) only at \(A\); its loss is \(3(0.065)+9(0.105)+9(0.070)=1.770\). Review reduces the forced minimum loss by \(0.1595\). Marking: review expectation (2), forced-cost rule/loss (1), MAP rule/loss and comparison (1).

\textit{Marking: use the displayed subpart allocations; award credit for the logical or numerical intermediate steps.}\par
{\footnotesize\textit{Source: PDF pp.11-12,15-18; supplementary asymmetric review-cost application.}}
\newpage\section*{Section D. Applications}
\textbf{30. Conditional independence, posterior decisions, and rejection}\hfill\textbf{[12 marks]}\par\medskip
\textbf{(a) [4 marks]} The joint-score pairs \((P(x,C_1),P(x,C_2))\) are
\[
00:(0.056,0.135),\quad01:(0.024,0.315),\quad
10:(0.224,0.045),\quad11:(0.096,0.105).
\]
For example, the \(C_1\) score at \(10\) is \(0.4(0.8)(0.7)=0.224\). The winning classes are \(C_2,C_2,C_1,C_2\), respectively. Summing each losing joint score gives the error
\[
0.056+0.024+0.045+0.096=0.221.
\]
Marking: four joint-score pairs (2), decisions (1), error sum (1).

\textbf{(b) [4 marks]} The feature-pair probabilities are \(0.191,0.339,0.269,0.201\). Their maximum posteriors are
\[
\frac{0.135}{0.191}\approx0.706806,\quad
\frac{0.315}{0.339}\approx0.929204,\quad
\frac{0.224}{0.269}\approx0.832714,\quad
\frac{0.105}{0.201}\approx0.522388.
\]
The acceptance threshold is \(1-0.30=0.70\). Accept \(00,01,10\) with the previous MAP labels; reject \(11\). Total risk is
\[
0.056+0.024+0.045+0.30(0.201)=0.1853.
\]
This includes review cost on all rejected inputs rather than treating them as free. Marking: posterior maxima/threshold (2), accepted/rejected pairs (1), expected risk (1).

\textbf{(c) [4 marks]} The marginal probabilities are
\[
P(X_1=1)=0.4(0.8)+0.6(0.25)=0.47,\quad
P(X_2=1)=0.4(0.3)+0.6(0.7)=0.54.
\]
Their joint probability is \(P(11)=0.096+0.105=0.201\), whereas \(0.47(0.54)=0.2538\). Since these differ, the features are dependent marginally. This is compatible with the stated conditional independence: both features respond to the shared, unobserved class variable. Marking: marginals (2), joint-versus-product comparison (1), conditional/marginal distinction (1).

\textit{Marking: use the displayed subpart allocations; award credit for the logical or numerical intermediate steps.}\par
{\footnotesize\textit{Source: PDF pp.11-13,16-18; supplementary marginal-dependence calculation.}}
\newpage\section*{Lecture Alignment and Coverage}
Primary source: \texttt{Bayesian\_Decision\_Theory.pdf} in \texttt{E:\textbackslash CS182\textbackslash 2026fall\textbackslash lecture}. All PDF pages count the cover as page 1. Additional arguments are marked as supplementary derivations when appropriate. These are original practice questions, not copied homework questions.

\textbf{T/F}
\begin{description}[leftmargin=12mm,itemsep=1pt]
\item[1] PDF pp.11-12.
\item[2] PDF p.15; supplementary numerical illustration.
\item[3] PDF p.23.
\item[4] PDF p.21; supplementary pointwise extension.
\item[5] PDF pp.6-7,16.
\item[6] PDF p.21.
\item[7] PDF pp.11-12,21.
\item[8] PDF pp.11-12,16.
\item[9] PDF p.15; supplementary randomization argument.
\item[10] PDF pp.12,16.
\item[11] PDF pp.17-18.
\item[12] PDF p.13.
\end{description}
\textbf{Multiple response}
\begin{description}[leftmargin=12mm,itemsep=1pt]
\item[13] PDF p.23; supplementary algebraic example.
\item[14] PDF pp.7,9-12,15-16; supplementary comparison of rule classes.
\item[15] PDF pp.17-18.
\item[16] PDF p.21; supplementary pointwise extension.
\item[17] PDF pp.15,20; supplementary loss transformations.
\item[18] PDF pp.11-12,21-22.
\item[19] PDF pp.11-13.
\item[20] PDF pp.11-12,16.
\item[21] PDF pp.15-18; supplementary tie/randomization interpretation.
\item[22] PDF pp.15-16; supplementary asymmetric-loss specialization.
\item[23] PDF p.15; supplementary randomization derivation.
\item[24] PDF pp.11-12,16.
\end{description}
\textbf{Proofs and applications}
\begin{description}[leftmargin=12mm,itemsep=1pt]
\item[25] PDF pp.15-16; supplementary integration and information comparison.
\item[26] PDF pp.15-18; supplementary asymmetric-cost extension.
\item[27] PDF pp.20-22; supplementary tie-preserving conditions.
\item[28] PDF p.23; supplementary affine and quadratic examples.
\item[29] PDF pp.11-12,15-18; supplementary asymmetric review-cost application.
\item[30] PDF pp.11-13,16-18; supplementary marginal-dependence calculation.
\end{description}
\end{document}
